\section{2015 Geometry \& Topology}

\subsection{Individual Exam}

\begin{exercise}
Let $n$, $m$ be positive integers. Show that the product of spheres $S^{n} \times S^{m}$ has trivial tangent bundle if and only if $n$ or $m$ is odd.
\end{exercise}

\begin{solution}
\textbf{Prerequisites:}
To understand this problem, one should be familiar with the following concepts:
1. \textbf{Tangent Bundle:} The collection of all tangent spaces of a manifold. A bundle is "trivial" (or the manifold is "parallelizable") if it is isomorphic to a product space $M \times \mathbb{R}^k$.
2. \textbf{Euler Characteristic ($\chi$):} A topological invariant. For a product of spaces, $\chi(X \times Y) = \chi(X)\chi(Y)$. For an $n$-sphere, $\chi(S^n) = 1 + (-1)^n$.
3. \textbf{Poincaré-Hopf Theorem:} Relates the Euler characteristic to the indices of vector fields. If a manifold has a trivial tangent bundle, it must have $\chi(M) = 0$.
4. \textbf{Stable Parallelizability:} A manifold $M$ is stably parallelizable if $TM \oplus \epsilon^1$ is trivial, where $\epsilon^1$ is a trivial line bundle. All spheres are stably parallelizable.
\par\vspace{\baselineskip}

\textbf{Steps:}
1. \textbf{The Necessity of $\chi = 0$:}
If $S^n \times S^m$ has a trivial tangent bundle, then it is parallelizable. A fundamental property of parallelizable manifolds is that they possess $k$ linearly independent non-vanishing vector fields (where $k$ is the dimension). By the Poincaré-Hopf theorem, if a manifold admits even a single non-vanishing vector field, its Euler characteristic must be zero.
\par\vspace{\baselineskip}

2. \textbf{Calculating the Euler Characteristic:}
The Euler characteristic of a sphere $S^k$ is given by $1 + (-1)^k$, which is $2$ if $k$ is even and $0$ if $k$ is odd. For the product $S^n \times S^m$, we have:
\[ \chi(S^n \times S^m) = \chi(S^n) \cdot \chi(S^m) = (1 + (-1)^n)(1 + (-1)^m). \]
For this product to be zero, at least one of the factors must be zero. This happens if and only if $n$ is odd or $m$ is odd. If both $n$ and $m$ are even, $\chi(S^n \times S^m) = 2 \times 2 = 4 \neq 0$.
\par\vspace{\baselineskip}

3. \textbf{The Sufficiency for Odd $n$ or $m$:}
Assume $n$ is odd. The sphere $S^n$ is known to have a non-vanishing vector field (the "hairy ball theorem" states this is only possible for odd dimensions). Furthermore, any sphere $S^n$ is "stably parallelizable" because its normal bundle in $\mathbb{R}^{n+1}$ is trivial, meaning $TS^n \oplus \epsilon^1$ is trivial.
\par\vspace{\baselineskip}

4. \textbf{Trivializing the Product Bundle:}
A known result in topology states that if a manifold $M$ is stably parallelizable and admits a non-vanishing vector field, then $M$ is parallelizable. Since $n$ is odd, $S^n$ has a non-vanishing vector field, and $S^n \times S^m$ is stably parallelizable (as both factors are). Thus, $S^n \times S^m$ is parallelizable, meaning its tangent bundle is trivial.
\par\vspace{\baselineskip}

\textbf{Intuition:}
Think of the tangent bundle as a way to assign a "frame" of directions to every point on a surface. On a circle ($S^1$), you can easily point "counter-clockwise" everywhere. On a 2-sphere ($S^2$), you always hit a "pole" where the directions get tangled (the Hairy Ball Theorem).
\par\vspace{\baselineskip}

When you multiply two spheres, you are creating a higher-dimensional surface. If at least one sphere is odd-dimensional (like a circle $S^1$ or $S^3$), it already has a "natural flow" that doesn't get stuck. This flow provides the "starting point" needed to straighten out the entire bundle of the combined space. If both spheres are even, the "tangles" from both spheres combine in a way that can never be fully combed smooth.
\end{solution}

\begin{exercise}
Show that there does not exist a compact three-dimensional manifold $M$ whose boundary is the real projective space $\mathbb{R}P^{2}$.
\end{exercise}

\begin{solution}
\textbf{Prerequisites:}
To approach this problem, one needs to be familiar with:
1. \textbf{Boundary of a Manifold:} A manifold $N$ is said to be a "boundary" if there exists a manifold $M$ such that $\partial M = N$.
2. \textbf{Cobordism Theory:} A branch of topology that studies manifolds by looking at what they can be the boundary of.
3. \textbf{Stiefel-Whitney Classes ($w_i$):} Topological invariants that measure the "twistedness" of a tangent bundle. They take values in the cohomology group $H^i(M; \mathbb{Z}_2)$.
4. \textbf{Stiefel-Whitney Numbers:} Evaluations of products of SW classes over the fundamental class of the manifold. A theorem by Pontryagin and Thom states that a manifold is a boundary if and only if all its Stiefel-Whitney numbers are zero.
\par\vspace{\baselineskip}

\textbf{Steps:}
1. \textbf{Defining the Characteristic Classes of $\mathbb{R}P^2$:}
The real projective space $\mathbb{R}P^2$ is the set of all lines through the origin in $\mathbb{R}^3$. Its total Stiefel-Whitney class $w(\mathbb{R}P^n)$ is given by the formula $(1+a)^{n+1}$, where $a$ is the generator of $H^1(\mathbb{R}P^n; \mathbb{Z}_2)$. For $n=2$, we have:
\[ w(\mathbb{R}P^2) = (1+a)^3 = 1 + 3a + 3a^2 \equiv 1 + a + a^2 \pmod 2. \]
This tells us that $w_1 = a$ and $w_2 = a^2$.
\par\vspace{\baselineskip}

2. \textbf{Calculating the Stiefel-Whitney Numbers:}
For a 2-dimensional manifold, the SW numbers are obtained by evaluating $w_1^2$ and $w_2$ on the fundamental class $[\mathbb{R}P^2]$.
- The first number is $w_1^2[\mathbb{R}P^2] = a^2[\mathbb{R}P^2]$. Since $a^2$ is the non-zero element in $H^2(\mathbb{R}P^2; \mathbb{Z}_2)$, this evaluation equals $1 \in \mathbb{Z}_2$.
- The second number is $w_2[\mathbb{R}P^2] = a^2[\mathbb{R}P^2] = 1$.
\par\vspace{\baselineskip}

3. \textbf{Applying the Thom Theorem:}
According to the theorem, for a manifold to be the boundary of a higher-dimensional manifold, all of its Stiefel-Whitney numbers must vanish. In our case:
\[ w_1^2[\mathbb{R}P^2] = 1 \neq 0. \]
Since we found at least one non-zero Stiefel-Whitney number, the manifold $\mathbb{R}P^2$ cannot be a boundary.
\par\vspace{\baselineskip}

\textbf{Intuition:}
Imagine trying to "fill in" a shape. A circle ($S^1$) is the boundary of a disk ($D^2$). A sphere ($S^2$) is the boundary of a ball ($D^3$). However, $\mathbb{R}P^2$ is "non-orientable" and has a strange topological twist (like a Möbius strip, but closed). 
\par\vspace{\baselineskip}

This "twist" is so fundamental that if you try to build a 3D object whose surface is $\mathbb{R}P^2$, you would eventually be forced to intersect the object with itself. There is no "interior" you can define that doesn't conflict with the "one-sided" nature of the surface. The Stiefel-Whitney numbers are just a rigorous mathematical way of detecting this impossible twist.
\end{solution}

\begin{exercise}
Let $M^{n}$ be a smooth manifold without boundary and $X$ a smooth vector field on $M$ which does not vanish at $p \in M$. Show that there exists a local coordinate chart $(U; x_{1}, \ldots, x_{n})$ centered at $p$ such that the vector field $X$ takes the form
\[ X = \frac{\partial}{\partial x_{1}}. \]
\end{exercise}

\begin{solution}
\textbf{Prerequisites:}
To solve this problem, one should understand:
1. \textbf{Vector Field:} An assignment of a vector to each point in a manifold.
2. \textbf{Flow of a Vector Field ($\Phi_t$):} The movement of points along the integral curves (trajectories) of the vector field. Mathematically, $\frac{d}{dt}\Phi_t(p) = X(\Phi_t(p))$.
3. \textbf{Hypersurface:} A sub-manifold of dimension $n-1$.
4. \textbf{Inverse Function Theorem:} States that if the derivative of a map is invertible at a point, the map is a local diffeomorphism (can be used as a coordinate chart).
\par\vspace{\baselineskip}

\textbf{Steps:}
1. \textbf{Selecting a Transversal Surface:}
Since $X(p) \neq 0$, we can choose a small $(n-1)$-dimensional hypersurface $S$ containing $p$ such that $X(p)$ is not tangent to $S$. We can assign local coordinates $(y_2, \dots, y_n)$ to $S$ such that $p$ corresponds to the origin.
\par\vspace{\baselineskip}

2. \textbf{Constructing the Map:}
Define a map $F$ from a neighborhood of the origin in $\mathbb{R}^n$ to $M$ by:
\[ F(t, y_2, \dots, y_n) = \Phi_t(q(y_2, \dots, y_n)), \]
where $q$ is the point on $S$ with coordinates $(y_2, \dots, y_n)$ and $\Phi_t$ is the flow of the vector field $X$.
\par\vspace{\baselineskip}

3. \textbf{Applying the Inverse Function Theorem:}
We compute the partial derivatives of $F$ at the origin $(0, \dots, 0)$:
- $\frac{\partial F}{\partial t} = \frac{d}{dt} \Phi_t(p) = X(p)$.
- $\frac{\partial F}{\partial y_i}$ are the tangent vectors to the surface $S$ at $p$.
Because we chose $S$ to be transversal to $X$, the vectors $\{X(p), \frac{\partial F}{\partial y_2}, \dots, \frac{\partial F}{\partial y_n}\}$ are linearly independent and form a basis for $T_p M$. Thus, $dF_0$ is invertible, and $F$ defines a local coordinate chart $(x_1, x_2, \dots, x_n)$ where $x_1 = t$ and $x_i = y_i$.
\par\vspace{\baselineskip}

4. \textbf{Verifying the Form of $X$:}
In these new coordinates, the vector field $X$ is given by the derivative of the flow with respect to time $t$. Since our first coordinate $x_1$ is exactly the time parameter $t$, we have:
\[ X = \frac{\partial}{\partial x_1}. \]
\par\vspace{\baselineskip}

\textbf{Intuition:}
Imagine a river flowing steadily. If you look at a small patch of water that isn't swirling (non-vanishing vector field), you can always imagine drawing a grid of lines: one set of lines follows the direction of the water's flow, and the others are perpendicular to it.
\par\vspace{\baselineskip}

This problem proves that you can always do this mathematically. By "straightening out" your coordinate grid so that the $x_1$-axis perfectly aligns with the direction of the flow, the complicated vector field $X$ becomes the simplest possible one: just a constant push in the $x_1$ direction. This is often called the "Flow Box" or "Rectification" theorem.
\end{solution}

\begin{exercise}
Let $M \subset \mathbb{R}^{3}$ be a compact simply-connected closed surface. Prove that if $M$ has constant mean curvature, then $M$ is a standard sphere.
\end{exercise}

\begin{solution}
\textbf{Prerequisites:}
Understanding this theorem requires knowledge of:
1. \textbf{Mean Curvature ($H$):} The average of the principal curvatures $k_1$ and $k_2$. $H = (k_1 + k_2)/2$.
2. \textbf{Umbilic Point:} A point on a surface where the two principal curvatures are equal ($k_1 = k_2$). At such points, the surface looks locally like a sphere.
3. \textbf{Second Fundamental Form:} A quadratic form that describes the local "bending" of a surface in $\mathbb{R}^3$.
4. \textbf{Hopf Differential:} A complex-valued quadratic differential $\Phi$ defined on the surface.
5. \textbf{Complex Structure on Surfaces:} Any simply-connected closed surface (genus 0) is conformally equivalent to the Riemann sphere $S^2$.
\par\vspace{\baselineskip}

\textbf{Steps:}
1. \textbf{Defining the Hopf Differential:}
On a surface, we can use a local complex coordinate $z$. The Hopf differential is defined as $\Phi = \phi(z) dz^2$, where $\phi(z) = \frac{1}{2}(h_{11} - h_{22}) - i h_{12}$, and $h_{ij}$ are the components of the second fundamental form.
\par\vspace{\baselineskip}

2. \textbf{Holomorphicity and Constant Mean Curvature:}
A key result in differential geometry is that if the mean curvature $H$ is constant, then the Hopf differential $\Phi$ is holomorphic (its derivative with respect to $\bar{z}$ is zero). This relates the extrinsic geometry (how the surface sits in space) to the intrinsic complex structure of the surface.
\par\vspace{\baselineskip}

3. \textbf{The Genus of the Surface:}
The problem states that $M$ is a simply-connected closed surface. By the classification of surfaces, this means $M$ is homeomorpic to $S^2$. A fundamental result from complex analysis states that there are no non-zero holomorphic quadratic differentials on a sphere ($S^2$ has genus $g=0$, and the space of such differentials has dimension 0).
\par\vspace{\baselineskip}

4. \textbf{Vanishing of the Hopf Differential:}
Since there are no non-zero holomorphic quadratic differentials on $S^2$, we must have $\Phi = 0$ everywhere on $M$. Looking back at the definition of $\Phi$, this implies $h_{11} = h_{22}$ and $h_{12} = 0$.
\par\vspace{\baselineskip}

5. \textbf{Conclusion from Umbilic Points:}
The condition $h_{11} = h_{22}$ and $h_{12} = 0$ means that every point on the surface is an umbilic point. A famous theorem in classical surface theory states that the only compact, connected surfaces in $\mathbb{R}^3$ where every point is umbilic are the round spheres.
\par\vspace{\baselineskip}

\textbf{Intuition:}
Imagine a soap bubble. The surface tension of the soap film tries to minimize the area, resulting in a constant mean curvature. If the bubble is "simple" (it doesn't have any holes or handles), it naturally forms a perfectly round shape.
\par\vspace{\baselineskip}

Mathematically, if the curvature were not constant or if there were "tensions" pulling in different directions, the surface would have to have "wrinkles" or points where it bends more in one direction than another (non-umbilic points). The fact that the surface is topologically a sphere and has constant "tension" (mean curvature) forces it to be perfectly symmetric in every direction at every point, which can only happen if it is a sphere.
\end{solution}

\begin{exercise}
Let $M$ be an $n$-dimensional compact Riemannian manifold with diameter $\pi/c$ and Ricci curvature $\geq (n-1)c^{2} > 0$. Show that $M$ is isometric to the standard $n$-sphere in $\mathbb{R}^{n+1}$ with radius $1/c$.
\end{exercise}

\begin{solution}
\textbf{Prerequisites:}
To understand this rigidity result, one should be familiar with:
1. \textbf{Ricci Curvature ($Ric$):} A measure of how the volume of a small ball in a manifold differs from the volume of a ball in Euclidean space. It is an average of sectional curvatures.
2. \textbf{Myers' Theorem:} States that if a complete Riemannian manifold has $Ric \ge (n-1)c^2 > 0$, then the manifold is compact and its diameter $D$ satisfies $D \leq \pi/c$.
3. \textbf{Bishop-Gromov Volume Comparison Theorem:} States that the volume of balls in a manifold with $Ric \ge (n-1)c^2$ grows no faster than the volume of balls in a sphere of constant curvature $c^2$.
4. \textbf{Isometry:} A map between two Riemannian manifolds that preserves the metric (distances and angles).
\par\vspace{\baselineskip}

\textbf{Steps:}
1. \textbf{The Boundary Case of Myers' Theorem:}
Myers' theorem provides an upper bound for the diameter of a manifold based on its Ricci curvature. In this problem, the diameter is given as exactly $\pi/c$. This means $M$ is the "largest" possible manifold allowed by its curvature. Such "boundary" cases in geometry often imply that the manifold has a very specific, rigid structure.
\par\vspace{\baselineskip}

2. \textbf{Using Volume Comparison:}
The Bishop-Gromov theorem tells us that the volume of $M$, denoted $V(M)$, is bounded by the volume of the standard $n$-sphere of radius $1/c$, denoted $V(S^n(1/c))$. Specifically, the ratio of volumes of balls in $M$ to balls in the sphere is a non-increasing function of the radius.
\par\vspace{\baselineskip}

3. \textbf{Diameter Rigidity:}
Cheng's Rigidity Theorem states that if the diameter of $M$ reaches the maximum value allowed by Myers' theorem ($D = \pi/c$), then the volume of $M$ must also be equal to the volume of the standard sphere $S^n(1/c)$.
\par\vspace{\baselineskip}

4. \textbf{Equality Case and Constant Curvature:}
In the Bishop-Gromov theorem, the equality $V(M) = V(S^n(1/c))$ holds if and only if the sectional curvatures of $M$ are constant and equal to $c^2$. This means $M$ is a space form.
\par\vspace{\baselineskip}

5. \textbf{Concluding Isometry:}
Since $M$ is a complete, compact manifold with constant sectional curvature $c^2$, it must be isometric to a quotient of the standard sphere. However, since its diameter is exactly $\pi/c$, it cannot be a smaller quotient (like $\mathbb{R}P^n$, which would have a smaller diameter). Thus, $M$ must be isometric to the standard $n$-sphere $S^n(1/c)$.
\par\vspace{\baselineskip}

\textbf{Intuition:}
Imagine you have a piece of rubber that you are trying to stretch out as far as possible. If the rubber is very "stiff" (has high positive curvature), you can't stretch it very far before it curves back on itself (Myers' theorem). 
\par\vspace{\baselineskip}

This problem says that if you managed to stretch that rubber to the absolute limit of what its stiffness allows, the only possible shape you could have created is a perfect, round sphere. Any other shape—any "bump" or "flat spot"—would have either made the manifold smaller or required less curvature. The "most stretched" manifold with a given minimum curvature is always the most symmetric one.
\end{solution}

\begin{exercise}
Suppose $(M, g)$ is a Riemannian manifold and $p \in M$. Show that the second-order Taylor series of $g$ in normal coordinates centered at $p$ is
\[ g_{ij}(x) = \delta_{ij} - \frac{1}{3}\sum_{k,l} R_{iklj}\,x_{k}x_{l} + O(|x|^{3}). \]
\end{exercise}

\begin{solution}
\textbf{Prerequisites:}
To tackle this problem, one must be familiar with:
1. \textbf{Normal Coordinates:} A local coordinate system centered at $p \in M$ where geodesics radiating from $p$ are straight lines through the origin. At $p=0$, the metric is Euclidean ($g_{ij}(0) = \delta_{ij}$).
2. \textbf{Christoffel Symbols ($\Gamma^k_{ij}$):} Functions that describe the connection (how vectors twist and turn as you move). In normal coordinates, $\Gamma^k_{ij}(0) = 0$.
3. \textbf{Riemann Curvature Tensor ($R_{ijkl}$):} A tensor that fundamentally measures the curvature of a manifold. It can be expressed in terms of second derivatives of the metric.
4. \textbf{Taylor Series in Multiple Variables:} The expansion of a function around a point using its partial derivatives.
\par\vspace{\baselineskip}

\textbf{Steps:}
1. \textbf{The Initial Taylor Expansion:}
We want to expand the metric tensor $g_{ij}(x)$ around $x=0$. The general formula is:
\[ g_{ij}(x) = g_{ij}(0) + \sum_k \partial_k g_{ij}(0) x_k + \frac{1}{2}\sum_{k,l} \partial_k \partial_l g_{ij}(0) x_k x_l + O(|x|^3). \]
By the definition of normal coordinates, $g_{ij}(0) = \delta_{ij}$. Also, because $\Gamma^k_{ij}(0) = 0$, all first partial derivatives of the metric vanish at the origin: $\partial_k g_{ij}(0) = 0$.
\par\vspace{\baselineskip}

2. \textbf{Relating Second Derivatives to Curvature:}
At the origin in normal coordinates, the Riemann curvature tensor simplifies to:
\[ R_{iklj}(0) = \frac{1}{2}(\partial_i \partial_l g_{kj} + \partial_k \partial_j g_{il} - \partial_i \partial_j g_{kl} - \partial_k \partial_l g_{ij}). \]
\par\vspace{\baselineskip}

3. \textbf{Using Symmetry and the Geodesic Equation:}
Because radial lines $x^i(t) = v^i t$ are geodesics, the geodesic equation $\ddot{x}^k + \Gamma^k_{ij}\dot{x}^i\dot{x}^j = 0$ implies that $\sum_{i,j} \Gamma^k_{ij}(x) x^i x^j = 0$ for all $x$. Differentiating this relation leads to the symmetry condition for the second derivatives of the metric at the origin:
\[ \partial_k \partial_l g_{ij}(0) + \partial_i \partial_l g_{jk}(0) + \partial_j \partial_l g_{ik}(0) = 0. \]
\par\vspace{\baselineskip}

4. \textbf{Solving for the Metric Term:}
By cyclically permuting the indices and using the symmetries of the Riemann tensor, we can combine the equations from Step 2 and Step 3. Specifically, one can derive the relation:
\[ \sum_{k,l} \partial_k \partial_l g_{ij}(0) x_k x_l = -\frac{2}{3} \sum_{k,l} R_{iklj} x_k x_l. \]
\par\vspace{\baselineskip}

5. \textbf{Final Substitution:}
Substituting this result back into the second-order term of our Taylor series (which had a factor of $\frac{1}{2}$), we get:
\[ g_{ij}(x) = \delta_{ij} - \frac{1}{3}\sum_{k,l} R_{iklj}x_kx_l + O(|x|^3). \]
\par\vspace{\baselineskip}

\textbf{Intuition:}
Imagine you are standing on a curved surface (like a sphere) and you lay down a flat piece of paper at your feet. Right where you are standing, the paper perfectly matches the surface (the metric is Euclidean, $\delta_{ij}$).
\par\vspace{\baselineskip}

As you move away from the center, the surface curves away from the flat paper. If you try to measure distances using the straight lines on your paper, your measurements will be increasingly wrong. This Taylor expansion tells you exactly how wrong you will be. The first-order "slope" is zero, but the second-order "bend" is determined precisely by the Riemann curvature tensor. The formula quantifies how the curvature forces the "true" distance to deviate from the "flat" approximation.
\end{solution}

\subsection{Team Exam}

\begin{exercise}
Let $SO(3)$ be the set of all $3 \times 3$ real matrices $A$ with determinant $1$ and satisfying ${}^{t}\!A\,A = I$, where $I$ is the identity matrix and ${}^{t}\!A$ is the transpose of $A$. Show that $SO(3)$ is a smooth manifold, and find its fundamental group. You need to prove your claims.
\end{exercise}

\begin{solution}
\textbf{Prerequisites:}
To solve this problem, you need to understand:
1. \textbf{Smooth Manifold:} A space that locally resembles Euclidean space $\mathbb{R}^n$ and allows for calculus.
2. \textbf{Regular Value Theorem:} A powerful tool to prove a set is a manifold. If $f: M \to N$ is a smooth map and $y \in N$ is a regular value (the differential $df$ is surjective everywhere on $f^{-1}(y)$), then $f^{-1}(y)$ is a smooth manifold.
3. \textbf{Fundamental Group ($\pi_1$):} A group formed by loops in a space, measuring its "holes."
4. \textbf{Covering Space:} A space that "unrolls" or "unwraps" another space. If a space has a simply connected cover, its fundamental group is isomorphic to the group of "deck transformations" (the symmetries of the unwrapping).
5. \textbf{Quaternions:} An extension of complex numbers. Unit quaternions form a 3-sphere $S^3$.
\par\vspace{\baselineskip}

\textbf{Steps:}
1. \textbf{Proving $SO(3)$ is a Manifold:}
Consider the space of all $3 \times 3$ real matrices, $M_{3 \times 3} \cong \mathbb{R}^9$. We define a map $f: M_{3 \times 3} \to \text{Sym}_{3 \times 3}$ (the space of symmetric $3 \times 3$ matrices, isomorphic to $\mathbb{R}^6$) by $f(A) = A^T A$. The orthogonal group $O(3)$ is the pre-image $f^{-1}(I)$.
\par\vspace{\baselineskip}

2. \textbf{Applying the Regular Value Theorem:}
We compute the differential $df_A(H) = A^T H + H^T A$. For any symmetric matrix $S \in \text{Sym}_{3 \times 3}$, we can find an $H$ such that $df_A(H) = S$ (specifically, $H = \frac{1}{2}A S$). This means $df_A$ is surjective at all $A \in f^{-1}(I)$. Therefore, $I$ is a regular value, making $O(3)$ a manifold of dimension $9 - 6 = 3$. Since $SO(3)$ is the connected component of $O(3)$ where $\det(A) = 1$, $SO(3)$ is also a smooth 3-dimensional manifold.
\par\vspace{\baselineskip}

3. \textbf{Finding the Fundamental Group:}
We use the unit quaternions, $S^3$. We can identify vectors in $\mathbb{R}^3$ with purely imaginary quaternions $v = xi + yj + zk$. For any unit quaternion $q \in S^3$, the operation $v \mapsto qvq^{-1}$ represents a rotation in $\mathbb{R}^3$. This gives us a continuous map $\rho: S^3 \to SO(3)$.
\par\vspace{\baselineskip}

4. \textbf{The Double Cover:}
The map $\rho$ is a surjective homomorphism. Furthermore, $q$ and $-q$ produce the exact same rotation, so the kernel of $\rho$ is $\{1, -1\}$. This means $\rho$ is a 2-to-1 covering map.
\par\vspace{\baselineskip}

5. \textbf{Conclusion for $\pi_1$:}
Because $S^3$ is a sphere of dimension 3, it is simply connected ($\pi_1(S^3) = 0$). It acts as the universal covering space of $SO(3)$. The fundamental group of $SO(3)$ is therefore isomorphic to the group of deck transformations of this cover, which consists of just two elements (the identity and the map $q \mapsto -q$). Thus, $\pi_1(SO(3)) \cong \mathbb{Z}_2$.
\par\vspace{\baselineskip}

\textbf{Intuition:}
Think about rotating an object in 3D space. You can describe any rotation by choosing an axis (a line through the object) and an angle to spin it.
\par\vspace{\baselineskip}

Imagine all these possible rotations forming a solid ball of radius $\pi$. The distance from the center is the angle of rotation, and the direction is the axis. A rotation of $\pi$ (180 degrees) is the maximum you ever need to go, because rotating further is just rotating less in the opposite direction. But wait! Rotating by $\pi$ clockwise around an axis is exactly the same physical state as rotating by $\pi$ counter-clockwise around the same axis. 
\par\vspace{\baselineskip}

This means opposite points on the surface of our "rotation ball" represent the exact same rotation. When you mathematically glue opposite points of a solid ball together, you get a strange space called the Real Projective Space ($\mathbb{R}P^3$). This space has a specific kind of "twist" in it, meaning if you take a certain path from the center out to the surface and back through the opposite side, you make a loop that can't be shrunk to a point unless you trace it twice. This is exactly what having a fundamental group of $\mathbb{Z}_2$ means: a "single twist" loop gets stuck, but a "double twist" loop can be untangled.
\end{solution}

\begin{exercise}
Let $X$ be a topological space. The suspension $S(X)$ of $X$ is the space obtained from $X \times [0, 1]$ by contracting $X \times \{0\}$ to a point and contracting $X \times \{1\}$ to another point. Describe the relation between the homology groups of $X$ and $S(X)$.
\end{exercise}

\begin{solution}
\textbf{Prerequisites:}
To understand this relation, you should know:
1. \textbf{Homology Groups ($H_k(X)$):} Algebraic structures that count the number of $k$-dimensional "holes" in a space.
2. \textbf{Reduced Homology ($\tilde{H}_k(X)$):} A slight variation of homology that ignores the "zero-dimensional holes" of a single connected component. It simplifies formulas.
3. \textbf{Cone over a Space ($CX$):} Formed by taking a cylinder $X \times [0,1]$ and crushing one end ($X \times \{1\}$) to a single point. A cone is always contractible (can be continuously shrunk to a point), so its reduced homology is zero.
4. \textbf{Mayer-Vietoris Sequence:} A mathematical tool that allows you to calculate the homology of a space by breaking it into two overlapping pieces, as long as you know the homology of the pieces and their intersection.
\par\vspace{\baselineskip}

\textbf{Steps:}
1. \textbf{Decomposing the Suspension:}
The suspension $S(X)$ can be viewed as the union of two open overlapping "cones." Let $C_1$ be the top cone formed by the points $X \times (1/3, 1]$ with $X \times \{1\}$ contracted. Let $C_2$ be the bottom cone formed by $X \times [0, 2/3)$ with $X \times \{0\}$ contracted.
\par\vspace{\baselineskip}

2. \textbf{Analyzing the Pieces:}
Both $C_1$ and $C_2$ are contractible spaces, which means their reduced homology groups are entirely trivial: $\tilde{H}_k(C_1) = \tilde{H}_k(C_2) = 0$ for all $k$.
\par\vspace{\baselineskip}

3. \textbf{Analyzing the Intersection:}
The intersection of the two cones, $C_1 \cap C_2$, is the "equator" band $X \times (1/3, 2/3)$. This cylinder is homotopically equivalent (can be continuously deformed) to the original space $X$. Therefore, $\tilde{H}_k(C_1 \cap C_2) \cong \tilde{H}_k(X)$.
\par\vspace{\baselineskip}

4. \textbf{Applying Mayer-Vietoris:}
We plug these pieces into the Mayer-Vietoris long exact sequence for reduced homology:
\[ \cdots \to \tilde{H}_k(C_1) \oplus \tilde{H}_k(C_2) \to \tilde{H}_k(S(X)) \xrightarrow{\partial} \tilde{H}_{k-1}(C_1 \cap C_2) \to \tilde{H}_{k-1}(C_1) \oplus \tilde{H}_{k-1}(C_2) \to \cdots \]
Substituting our known values (the zeros and $X$), this segment becomes:
\[ 0 \to \tilde{H}_k(S(X)) \xrightarrow{\partial} \tilde{H}_{k-1}(X) \to 0 \]
\par\vspace{\baselineskip}

5. \textbf{The Isomorphism:}
Because the sequence is exact and is bracketed by zeros, the map in the middle must be an isomorphism. This gives us the fundamental relation:
\[ \tilde{H}_k(S(X)) \cong \tilde{H}_{k-1}(X) \text{ for all } k. \]
\par\vspace{\baselineskip}

\textbf{Intuition:}
Imagine taking a simple rubber band (a 1-dimensional circle, $S^1$). To suspend it, you pull the top edge up to a single point (the North Pole) and pull the bottom edge down to a single point (the South Pole). You have just created a hollow balloon (a 2-dimensional sphere, $S^2$).
\par\vspace{\baselineskip}

The suspension operation fundamentally "stretches" a space into the next higher dimension. A point becomes a line segment, a circle becomes a sphere, a 2D sphere becomes a 3D sphere, etc. Naturally, whatever "holes" existed in your original shape get stretched as well. A 1D hole (a loop) in $X$ gets stretched up and down to become a 2D hole (a hollow cavity) in $S(X)$. Therefore, the $k$-dimensional homology of the suspended space is exactly the $(k-1)$-dimensional homology of the original space.
\end{solution}

\begin{exercise}
Let $F : M \to N$ be a smooth map between two manifolds. Let $X_{1}, X_{2}$ be smooth vector fields on $M$ and let $Y_{1}, Y_{2}$ be smooth vector fields on $N$. Prove that if $Y_{1} = F_{*}X_{1}$ and $Y_{2} = F_{*}X_{2}$, then $F_{*}[X_{1}, X_{2}] = [Y_{1}, Y_{2}]$, where $[~, ~]$ is the Lie bracket.
\end{exercise}

\begin{solution}
\textbf{Prerequisites:}
To understand this proof, one needs to be familiar with:
1. \textbf{Vector Fields as Derivations:} A vector field $X$ on a manifold acts on a smooth function $f$ to produce a new function $X(f)$ (the directional derivative of $f$ along $X$).
2. \textbf{Pushforward ($F_*$):} A map $F: M \to N$ "pushes" tangent vectors from $M$ to $N$. If a vector field $X$ on $M$ pushes forward to $Y$ on $N$ ($Y = F_*X$), it means that for any function $g$ on $N$, taking the derivative along $X$ on $M$ is the same as taking the derivative along $Y$ on $N$. Mathematically: $X(g \circ F) = (Yg) \circ F$.
3. \textbf{Lie Bracket ($[X, Y]$):} An operation that combines two vector fields to create a third one, defined by its action on a function $f$: $[X, Y](f) = X(Y(f)) - Y(X(f))$.
\par\vspace{\baselineskip}

\textbf{Steps:}
1. \textbf{Setting up the Goal:}
To prove that $F_*[X_1, X_2] = [Y_1, Y_2]$, we must show that for any smooth "test function" $g$ defined on $N$, the action of both sides on $g$ yields the same result. That is, we must show:
\[ [X_1, X_2](g \circ F) = ([Y_1, Y_2]g) \circ F. \]
\par\vspace{\baselineskip}

2. \textbf{Expanding the Left Side:}
Using the definition of the Lie bracket on the manifold $M$, we expand the left side acting on the composite function $(g \circ F)$:
\[ [X_1, X_2](g \circ F) = X_1(X_2(g \circ F)) - X_2(X_1(g \circ F)). \]
\par\vspace{\baselineskip}

3. \textbf{Applying the Pushforward Condition (First Time):}
We are given that $Y_2 = F_*X_2$, which means $X_2(g \circ F) = (Y_2 g) \circ F$. Similarly, $X_1(g \circ F) = (Y_1 g) \circ F$. Substituting these into our expanded equation gives:
\[ = X_1((Y_2 g) \circ F) - X_2((Y_1 g) \circ F). \]
\par\vspace{\baselineskip}

4. \textbf{Applying the Pushforward Condition (Second Time):}
Now, treat $(Y_2 g)$ and $(Y_1 g)$ as new functions defined on $N$. We apply the pushforward conditions ($Y_1 = F_*X_1$ and $Y_2 = F_*X_2$) again:
\[ = (Y_1(Y_2 g)) \circ F - (Y_2(Y_1 g)) \circ F. \]
\par\vspace{\baselineskip}

5. \textbf{Factoring and Concluding:}
We can factor out the composition with $F$ on the right side:
\[ = (Y_1 Y_2 g - Y_2 Y_1 g) \circ F. \]
By the definition of the Lie bracket on the manifold $N$, this is exactly:
\[ = ([Y_1, Y_2]g) \circ F. \]
Since this holds for any function $g$, the vector fields must be equal: $F_*[X_1, X_2] = [Y_1, Y_2]$.
\par\vspace{\baselineskip}

\textbf{Intuition:}
Imagine two rivers (vector fields $X_1$ and $X_2$) flowing across a landscape $M$. The Lie bracket $[X_1, X_2]$ measures how much a boat would be "knocked off course" if it tried to float down $X_1$ for a minute, then $X_2$ for a minute, compared to doing it in the reverse order.
\par\vspace{\baselineskip}

Now, suppose you have a magical camera (the map $F$) that projects an image of landscape $M$ onto a screen $N$. The pushforward $F_*$ is just the projection of the rivers onto the screen (becoming rivers $Y_1$ and $Y_2$). 
\par\vspace{\baselineskip}

This proof simply says: if you watch the boat being knocked off course on the real landscape ($[X_1, X_2]$), and then take a picture of it ($F_*$), it looks exactly the same as if you just watched the image of the boat being knocked off course by the projected rivers ($[Y_1, Y_2]$) on the screen. The "failure to commute" is perfectly preserved by the projection.
\end{solution}

\begin{exercise}
Let $M_{1}$ and $M_{2}$ be two compact convex closed surfaces in $\mathbb{R}^{3}$, and $f : M_{1} \to M_{2}$ a diffeomorphism such that $M_{1}$ and $M_{2}$ have the same inner normal vectors and Gauss curvatures at the corresponding points. Prove that $f$ is a translation.
\end{exercise}

\begin{solution}
\textbf{Prerequisites:}
To approach this problem, one must understand:
1. \textbf{Gauss Map ($N$):} A map from a surface $M$ to the unit sphere $S^2$ that takes a point on the surface to its unit normal vector. For a strictly convex closed surface, this map is a diffeomorphism (a smooth 1-to-1 correspondence).
2. \textbf{Gauss Curvature ($K$):} A measure of the intrinsic curvature of a surface. For convex surfaces, $K > 0$.
3. \textbf{Support Function ($h$):} A function $h: S^2 \to \mathbb{R}$ defined as $h(u) = \langle x(u), u \rangle$, where $x(u)$ is the point on the surface with normal $u$. It measures the distance from the origin to the tangent plane at $x(u)$.
4. \textbf{Maximum Principle for Elliptic PDEs:} A mathematical tool showing that solutions to certain differential equations cannot have "local bumps" without being constant or having a specific rigid form.
\par\vspace{\baselineskip}

\textbf{Steps:}
1. \textbf{Parameterizing over the Sphere:}
Because both $M_1$ and $M_2$ are compact and convex, we can parameterize both surfaces using the inverse Gauss map. A point $u \in S^2$ corresponds to exactly one point $x_1(u)$ on $M_1$ and one point $x_2(u)$ on $M_2$. The diffeomorphism $f$ essentially just maps $x_1(u)$ to $x_2(u)$.
\par\vspace{\baselineskip}

2. \textbf{The Support Function Equation:}
For any convex surface, the principal radii of curvature (the inverses of the principal curvatures) are the eigenvalues of the matrix $(h_{ij} + h \delta_{ij})$ on $S^2$, where $h_{ij}$ is the covariant Hessian of the support function $h$.
\par\vspace{\baselineskip}

3. \textbf{Relating to Gauss Curvature:}
The product of the principal radii of curvature is $1/K$. Therefore, the support function $h$ must satisfy the Monge-Ampère equation:
\[ \det(h_{ij} + h \delta_{ij}) = \frac{1}{K(u)}. \]
\par\vspace{\baselineskip}

4. \textbf{Applying the Maximum Principle:}
We are given that $K_1(u) = K_2(u)$ for all directions $u$. Therefore, both $h_1$ and $h_2$ are solutions to the exact same highly non-linear differential equation on the sphere. Let $w = h_1 - h_2$. By applying the maximum principle to the difference of the two Monge-Ampère equations, one can show that $w$ must be an eigenfunction of the Laplacian on the sphere corresponding to the eigenvalue $-2$.
\par\vspace{\baselineskip}

5. \textbf{Concluding the Translation:}
The eigenfunctions of the Laplacian on $S^2$ for eigenvalue $-2$ are exactly the linear functions of the coordinates in $\mathbb{R}^3$. This means $w(u) = \langle a, u \rangle$ for some constant vector $a \in \mathbb{R}^3$.
\par\vspace{\baselineskip}

6. \textbf{Final Geometric Implication:}
Since $h_1(u) - h_2(u) = \langle a, u \rangle$, the distance to the tangent planes differs exactly by the projection of a constant vector $a$. This implies that the position vectors themselves satisfy $x_1(u) = x_2(u) + a$. Therefore, $M_1$ is just $M_2$ shifted by the vector $a$ (a translation).
\par\vspace{\baselineskip}

\textbf{Intuition:}
Imagine you have an invisible solid object. I tell you exactly how "pointy" or "flat" the object is (the Gauss curvature) at every single possible angle you could look at it from (the normal vector). 
\par\vspace{\baselineskip}

Minkowski's uniqueness theorem (which this problem proves) states that this information completely dictates the shape of the object. If you and your friend both have an object that matches my description perfectly, you must have the exact same shape. The only difference is that you might be holding your object in a different location in the room than your friend is. No matter how you try to "dent" or "warp" your object while keeping it convex, changing the shape will inevitably change the curvature at some orientation.
\end{solution}

\begin{exercise}
Prove the second Bianchi identity:
\[ R_{ijkl,h} + R_{ijlh,k} + R_{ijhk,l} = 0. \]
\end{exercise}

\begin{solution}
\textbf{Prerequisites:}
To prove this identity, you need to be familiar with:
1. \textbf{Riemann Curvature Tensor ($R_{ijkl}$):} The fundamental object that measures curvature in a Riemannian manifold.
2. \textbf{Covariant Derivative (denoted by a comma, e.g., $R_{ijkl,h}$):} The generalization of the partial derivative for tensors, which accounts for the "twisting" of the coordinate system (involving Christoffel symbols $\Gamma$).
3. \textbf{Normal Coordinates:} A local coordinate system around a point $p$ where the Christoffel symbols vanish ($\Gamma^m_{ij}(p) = 0$). This drastically simplifies covariant derivatives at $p$.
4. \textbf{Tensor Nature:} If an equation made entirely of tensors is true in one coordinate system, it is true in all coordinate systems.
\par\vspace{\baselineskip}

\textbf{Steps:}
1. \textbf{Choosing the Right Coordinates:}
Because the Second Bianchi Identity is a tensor equation (both sides are tensors), it is sufficient to prove it at a single arbitrary point $p \in M$ using normal coordinates centered at $p$.
\par\vspace{\baselineskip}

2. \textbf{Simplifying the Covariant Derivative:}
In normal coordinates at $p$, all Christoffel symbols $\Gamma$ are zero. Because the covariant derivative is just the standard partial derivative plus a bunch of terms involving $\Gamma$, at point $p$, the covariant derivative simply becomes the regular partial derivative:
\[ R_{ijkl,h}(p) = \partial_h R_{ijkl}(p). \]
\par\vspace{\baselineskip}

3. \textbf{Expanding the Riemann Tensor:}
We recall the formula for the Riemann tensor in terms of the metric $g$. At the origin of normal coordinates, it is given entirely by second partial derivatives of the metric:
\[ R_{ijkl} = \frac{1}{2} (\partial_i \partial_l g_{jk} + \partial_j \partial_k g_{il} - \partial_i \partial_k g_{jl} - \partial_j \partial_l g_{ik}). \]
\par\vspace{\baselineskip}

4. \textbf{Taking the Derivative:}
Now we take the partial derivative $\partial_h$ of the above expression. This gives us terms with third partial derivatives of the metric:
\[ R_{ijkl,h} = \partial_h R_{ijkl} = \frac{1}{2} (\partial_h \partial_i \partial_l g_{jk} + \partial_h \partial_j \partial_k g_{il} - \partial_h \partial_i \partial_k g_{jl} - \partial_h \partial_j \partial_l g_{ik}). \]
\par\vspace{\baselineskip}

5. \textbf{Cyclic Permutation and Cancellation:}
The Second Bianchi Identity asks us to cyclically permute the last three indices $(k, l, h)$ and sum the results:
\[ R_{ijkl,h} + R_{ijlh,k} + R_{ijhk,l}. \]
We write out all three expressions using the formula from Step 4. Because partial derivatives commute (e.g., $\partial_h \partial_i \partial_l g_{jk} = \partial_l \partial_i \partial_h g_{jk}$), we can rearrange the order of derivatives. When we add the three massive expressions together, every single term appears exactly twice: once with a positive sign and once with a negative sign. They all perfectly cancel out, leaving exactly $0$.
\par\vspace{\baselineskip}

6. \textbf{Conclusion:}
Since the tensor equation holds in normal coordinates at point $p$, and $p$ was arbitrary, the identity holds globally across the entire manifold.
\par\vspace{\baselineskip}

\textbf{Intuition:}
In basic calculus, we learn that "the derivative of a constant is zero." In topology, there is a fundamental principle that "the boundary of a boundary is zero" (a 2D disk has a 1D circle as a boundary, but a 1D circle has no endpoints, so its boundary is 0).
\par\vspace{\baselineskip}

The Second Bianchi Identity is the differential geometry version of this topological rule. The Riemann tensor measures how much a vector changes when you move it around a small closed loop (a boundary). Taking the covariant derivative of the Riemann tensor ($R_{ijkl,h}$) is akin to looking at how this change itself varies as you shift the loop. Summing the cyclical permutations is mathematically equivalent to taking the "boundary of the boundary." The identity guarantees that the geometry remains mathematically consistent and "closed" without tearing. In Einstein's General Relativity, this exact identity is what ensures the conservation of energy and momentum in the universe!
\end{solution}

\begin{exercise}
Let $M_{1}, M_{2}$ be two complete $n$-dimensional Riemannian manifolds and $\gamma_{i} : [0, a] \to M_{i}$ are two arc length parametrized geodesics. Let $\rho_{i}$ be the distance function to $\gamma_{i}(0)$ on $M_{i}$. Assume that $\gamma_{i}(a)$ is within the cut locus of $\gamma_{i}(0)$ and for any $0 \leq t \leq a$ we have the inequality of sectional curvatures
\[ K_{1}\!\left(X_{1}, \frac{\partial}{\partial \gamma_{1}}\right) \geq K_{2}\!\left(X_{2}, \frac{\partial}{\partial \gamma_{2}}\right), \]
where $X_{i} \in T_{\gamma_{i}(t)}M_{i}$ is any unit vector orthogonal to the tangent $\frac{\partial}{\partial \gamma_{i}}$.
Then
\[ \operatorname{Hess}(\rho_{1})(\tilde{X}_{1}, \tilde{X}_{1}) \leq \operatorname{Hess}(\rho_{2})(\tilde{X}_{2}, \tilde{X}_{2}), \]
where $\tilde{X}_{i} \in T_{\gamma_{i}(a)}M_{i}$ is any unit vector orthogonal to the tangent $\frac{\dot{\sigma}}{\partial \gamma_{i}}(a)$.
\end{exercise}

\begin{solution}
\textbf{Prerequisites:}
To grasp this comparison theorem, one needs to know:
1. \textbf{Sectional Curvature ($K$):} A measure of the curvature of a 2D slice of the manifold.
2. \textbf{Distance Function ($\rho$):} The function giving the shortest distance from a starting point $p = \gamma(0)$.
3. \textbf{Hessian ($\operatorname{Hess}$):} The second derivative of a function. The Hessian of the distance function measures how quickly the "distance spheres" (sets of points at a constant distance from $p$) are bending.
4. \textbf{Cut Locus:} The point where geodesics from $p$ start to intersect or fail to be the shortest path. Before the cut locus, distance is smooth.
5. \textbf{Riccati Equation:} A specific type of non-linear differential equation that frequently appears in geometry when describing how things spread out or converge.
\par\vspace{\baselineskip}

\textbf{Steps:}
1. \textbf{The Shape Operator of Distance Spheres:}
Let $U(t)$ be the matrix representing the shape operator (the Hessian of the distance function $\rho$) of the distance sphere at radius $t$ along the geodesic $\gamma$. By definition, $\operatorname{Hess}(\rho)(X, X) = \langle U(t)X, X \rangle$.
\par\vspace{\baselineskip}

2. \textbf{The Matrix Riccati Equation:}
As we move along the geodesic $\gamma(t)$, the rate at which the distance spheres bend ($U(t)$) is governed by the curvature of the space. This is beautifully captured by the matrix Riccati equation:
\[ U'(t) + U(t)^2 + R(t) = 0, \]
where $R(t)$ is the curvature operator defined by $R(X) = R(X, \dot{\gamma})\dot{\gamma}$.
\par\vspace{\baselineskip}

3. \textbf{Comparing Curvatures:}
The problem states that the sectional curvature of $M_1$ is everywhere greater than or equal to the sectional curvature of $M_2$ in the direction of the geodesics. Mathematically, this translates to the inequality of the curvature operators: $R_1(t) \ge R_2(t)$ (meaning the difference is a positive semi-definite matrix).
\par\vspace{\baselineskip}

4. \textbf{The Sturm Comparison Theorem:}
Because both geodesics start at $t=0$, near the origin, the distance spheres look like standard Euclidean spheres, meaning $U_i(t) \sim \frac{1}{t} I$ as $t \to 0$. We have two Riccati equations with the same initial behavior, but the first one has a larger "curvature penalty" term ($R_1 \ge R_2$). The Sturm comparison theorem for Riccati equations guarantees that the solution with the larger penalty term gets pushed down faster. Therefore, for all $t \in (0, a]$ before the cut locus:
\[ U_1(t) \le U_2(t). \]
\par\vspace{\baselineskip}

5. \textbf{Concluding the Hessian Inequality:}
Since $U_1 \le U_2$ as matrices, multiplying them by unit vectors yields the final result at $t=a$:
\[ \operatorname{Hess}(\rho_1)(\tilde{X}_1, \tilde{X}_1) \le \operatorname{Hess}(\rho_2)(\tilde{X}_2, \tilde{X}_2). \]
\par\vspace{\baselineskip}

\textbf{Intuition:}
Imagine tossing a pebble into a pond. The ripples (distance spheres) spread out evenly. Now, imagine doing this on two different surfaces: a flat table ($M_2$, zero curvature) and the top of a round bowling ball ($M_1$, positive curvature).
\par\vspace{\baselineskip}

As the ripple spreads on the bowling ball, the positive curvature of the ball causes the surface to "bend away" from the straight paths. This causes the geodesics (the paths the water takes) to converge towards each other much faster than they would on the flat table.
\par\vspace{\baselineskip}

The Hessian of the distance function measures precisely how "tightly" these distance spheres are curving. Because the higher curvature of $M_1$ forces the geodesics to converge, the distance spheres on $M_1$ become more sharply bent (or shrink faster) compared to $M_2$. In geometric terms, a "sharper inward bend" corresponds to a mathematically smaller (more negative) Hessian. Thus, higher curvature strictly leads to a smaller Hessian.
\end{solution}
